\section{Space Complexity --- Machines --- Bin}
%%% generated by the blueprint pipeline from TCSlib.Complexity.SpaceComplexity.Machines.Bin
\subsection*{Space Complexity --- Machines --- Bin: Introduction}

Declarations of the Lean module \texttt{TCSlib.Complexity.SpaceComplexity.Machines.Bin}.

\subsection*{Space Complexity --- Machines --- Bin: Results}

\begin{definition}[Binary value of a Boolean list]
\lean{Complexity.LogProg.bitsVal}
\label{Complexity.LogProg.bitsVal}
\leanok
For a finite list of Booleans, define its value recursively by assigning the empty list
value $0$ and assigning a list with first bit $b$ and remaining list $w$ the value
$b+2\,\mathrm{value}(w)$, where false has value $0$ and true has value $1$.
\end{definition}

\begin{definition}[Incrementing a binary list]
\lean{Complexity.LogProg.incW}
\label{Complexity.LogProg.incW}
\leanok
Define an operation on finite Boolean lists by sending the empty list to the one-bit
list containing true; a list beginning with false is changed to begin with true, and a
list beginning with true is changed to begin with false while the remaining list is
recursively incremented.
\end{definition}

\begin{lemma}[Bits of a doubled number plus a bit]
\lean{Complexity.LogProg.bits_two_mul_add}
\label{Complexity.LogProg.bits_two_mul_add}
\leanok
\difficulty{4}
For every natural number $n$ and bit $b$, if $n=0$ implies $b=1$, then the binary bit
list of $2n+b$ is the list formed by placing $b$ before the binary bit list of $n$.
\end{lemma}

\begin{lemma}[Value of the binary expansion word]
\lean{Complexity.LogProg.bitsVal_bits}
\label{Complexity.LogProg.bitsVal_bits}
\leanok
\uses{Complexity.LogProg.bitsVal, Complexity.LogProg.bits_two_mul_add}
\difficulty{6}
For every natural number $n$, the value of the binary word associated to $n$ is $n$.
\end{lemma}

\begin{lemma}[Injectivity of binary expansion]
\lean{Complexity.LogProg.bits_injective}
\label{Complexity.LogProg.bits_injective}
\leanok
\uses{Complexity.LogProg.bitsVal_bits}
\difficulty{1}
Distinct natural numbers have distinct binary bit lists.
\end{lemma}

\begin{lemma}[Binary expansion of a successor]
\lean{Complexity.LogProg.bits_succ}
\label{Complexity.LogProg.bits_succ}
\leanok
\uses{Complexity.LogProg.bits_two_mul_add, Complexity.LogProg.incW}
\difficulty{5}
For every natural number $n$, the binary bit list of $n+1$ is obtained by changing each
initial $1$ to $0$ until the first $0$, which is changed to $1$; if the list has no such
$0$, a $1$ is appended after changing all its bits to $0$.
\end{lemma}

\begin{lemma}[Length Bound for Binary Expansion]
\lean{Complexity.LogProg.length_bits_le}
\label{Complexity.LogProg.length_bits_le}
\leanok
\difficulty{1}
For natural numbers $n$ and $m$, if $n<2^m$, then the binary expansion of $n$ has length
at most $m$.
\end{lemma}

\begin{lemma}[Final bit of a nonempty binary expansion]
\lean{Complexity.LogProg.bits_getLast}
\label{Complexity.LogProg.bits_getLast}
\leanok
\uses{Complexity.LogProg.bitsVal, Complexity.LogProg.bitsVal_bits, Complexity.LogProg.bits_two_mul_add}
\difficulty{6}
For every natural number $n$, if its binary bit list is nonempty, then the last bit of
that list is $1$.
\end{lemma}

\begin{lemma}[Logarithmic bound for a polynomial expression]
\lean{Complexity.LogProg.log_poly_bound}
\label{Complexity.LogProg.log_poly_bound}
\leanok
\uses{Complexity.LogProg.succ_le_two_pow_log}
\difficulty{5}
For all nonnegative integers $A,c,B$, there is a nonnegative integer $K$ such that, for
every nonnegative integer $n$, the base-two logarithm of $A(n+1)^c+B$, rounded down and
then increased by one, is at most $K$ times the base-two logarithm of $n$, rounded down
and then increased by one.
\end{lemma}

%%% added by the blueprint pipeline
\subsection*{Space Complexity --- Machines --- Bin: Further results}

\begin{lemma}[Bound on the value of a bit string]
\lean{Complexity.LogProg.bitsVal_lt_two_pow}
\label{Complexity.LogProg.bitsVal_lt_two_pow}
\leanok
\uses{Complexity.LogProg.bitsVal}
\difficulty{4}
For every finite bit string $w$, its value is strictly less than $2^{|w|}$, where $|w|$
is its length.
\end{lemma}

\begin{lemma}[Uniqueness of Fixed-Length Binary Words]
\lean{Complexity.LogProg.eq_of_bitsVal_eq}
\label{Complexity.LogProg.eq_of_bitsVal_eq}
\leanok
\uses{Complexity.LogProg.bitsVal}
\difficulty{4}
If two finite Boolean words have the same length and the same value under the binary
evaluation that assigns the head bit weight $1$ and doubles the value of the remaining
suffix, then the words are equal.
\end{lemma}

\begin{lemma}[Monotonicity of Binary Expansion Length]
\lean{Complexity.LogProg.length_bits_mono}
\label{Complexity.LogProg.length_bits_mono}
\leanok
\difficulty{3}
If natural numbers $a$ and $b$ satisfy $a\le b$, then the length of the binary expansion
of $a$ is at most the length of the binary expansion of $b$.
\end{lemma}

\begin{lemma}[Logarithmic Bound for Binary Expansion Length]
\lean{Complexity.LogProg.length_bits_le_log}
\label{Complexity.LogProg.length_bits_le_log}
\leanok
\uses{Complexity.LogProg.length_bits_le}
\difficulty{1}
For every natural number $n$, the length of its binary expansion is at most
$\lfloor\log_2 n\rfloor+1$.
\end{lemma}

\begin{lemma}[A bound from the binary logarithm]
\lean{Complexity.LogProg.succ_le_two_pow_log}
\label{Complexity.LogProg.succ_le_two_pow_log}
\leanok
\difficulty{1}
For every natural number $n$, if $k$ is the greatest natural number such that $2^k\le n$
(with $k=0$ when $n=0$), then $n+1\le 2^{k+1}$.
\end{lemma}
